% --- UNIVERSAL PREAMBLE BLOCK FOR BEAMER ---
\documentclass[11pt, aspectratio=169]{beamer}
\usepackage{fontspec}

\usepackage[portuguese, bidi=basic, provide=*]{babel}

\babelprovide[import, onchar=ids fonts]{portuguese}
\babelprovide[import, onchar=ids fonts]{english}

% Beamer uses sans-serif fonts by default, so we set Noto Sans for both rm and sf
\babelfont{rm}{Noto Sans}
\babelfont{sf}{Noto Sans}

% Theme selection
\usetheme{Madrid}
\usecolortheme{whale}

% Standard math packages
\usepackage{amsmath}
\usepackage{amssymb}
\usepackage{booktabs}

% Title information
\title[Método do Potencial]{Seminário 1: Método do Potencial em Jogos Posicionais}
\subtitle{Análise de Erdős e Selfridge (1973)}
\author{Equipe 1}
\institute{Combinatória Extremal e Probabilística}
\date{}

\begin{document}

\begin{frame}
    \titlepage
\end{frame}

\begin{frame}{Objetivos do Seminário}
    \begin{itemize}
        \item \textbf{Compreender a dinâmica} de um jogo posicional do tipo Maker-Breaker.
        \item \textbf{Apresentar a Técnica do Potencial:} Como uma função matemática pode substituir a análise exaustiva da árvore de jogadas.
        \item \textbf{Demonstrar o Critério de Erdős-Selfridge:} Um teorema clássico e global que garante a vitória/empate do jogador defensivo.
        \item \textbf{Discutir a Otimalidade:} Analisar a construção extremal exata para o caso uniforme.
        \item \textbf{Aplicação Prática:} O jogo de cliques em grafos ($K_N \to K_q$).
    \end{itemize}
\end{frame}

\begin{frame}{O que são Jogos Maker-Breaker?}
    \begin{block}{Definição Formal}
        Um jogo Maker-Breaker é definido por um par $(X, \mathcal{F})$, onde:
        \begin{itemize}
            \item $X$ é um conjunto finito (o \textbf{tabuleiro}).
            \item $\mathcal{F} \subseteq 2^X$ é uma família de subconjuntos de $X$ (os \textbf{conjuntos vencedores}).
        \end{itemize}
    \end{block}
    
    \begin{block}{Regras e Dinâmica}
        \begin{itemize}
            \item \textbf{Maker} e \textbf{Breaker} alternam turnos escolhendo um elemento livre $x \in X$.
            \item O \textbf{Maker} vence se conseguir ocupar integralmente algum conjunto $A \in \mathcal{F}$.
            \item O \textbf{Breaker} vence se conseguir colocar pelo menos uma de suas marcas em TODOS os conjuntos de $\mathcal{F}$, impedindo o Maker de vencer.
        \end{itemize}
    \end{block}
\end{frame}

\begin{frame}{O Critério de Erdős-Selfridge (1973)}
    \begin{block}{Teorema Principal}
        Em um jogo Maker-Breaker sobre o hipergrafo $\mathcal{F}$, se o Maker joga primeiro e a seguinte desigualdade estrita é satisfeita:
        \begin{equation*}
            \sum_{A \in \mathcal{F}} 2^{-|A|} < \frac{1}{2}
        \end{equation*}
        então o \textbf{Breaker possui uma estratégia vencedora} (ele consegue impedir a vitória do Maker, independentemente de quão bem o Maker jogue).
    \end{block}
    \vspace{0.5cm}
    \textit{Nota: Se o Breaker fosse o primeiro a jogar, o lado direito da desigualdade seria $1$.}
\end{frame}

\begin{frame}{Técnica A: A Função Potencial}
    Para provar o teorema, Erdős e Selfridge introduziram uma \textbf{função potencial} (um método "quase-probabilístico") para avaliar o estado do jogo.
    
    \begin{block}{Definindo o Peso (Potencial)}
        Após a rodada $i$, sejam $M_i$ as peças do Maker e $B_i$ as do Breaker.
        Um conjunto $A \in \mathcal{F}$ está \textbf{ativo} (ou sobrevivente) se $A \cap B_i = \emptyset$.
        O potencial do tabuleiro é a soma dos pesos dos conjuntos ativos:
        \begin{equation*}
            \Phi_i = \sum_{\substack{A \in \mathcal{F} \\ A \cap B_i = \emptyset}} 2^{-|A \setminus M_i|}
        \end{equation*}
    \end{block}
    
    \begin{itemize}
        \item O expoente $|A \setminus M_i|$ é o número de peças que \textbf{faltam} para o Maker completar $A$.
        \item Quanto mais perto o Maker está de completar o conjunto, maior o seu peso.
    \end{itemize}
\end{frame}

\begin{frame}{Dinâmica das Jogadas: O Turno do Maker}
    \begin{itemize}
        \item No início do jogo (rodada 0), o potencial é estritamente menor que $1/2$ (por hipótese): 
        $$\Phi_0 = \sum_{A \in \mathcal{F}} 2^{-|A|} < 1/2$$
        \item \textbf{Ação do Maker:} O Maker escolhe um elemento $x$.
        \item Para todo conjunto ativo $A$ que contém $x$, o número de elementos faltantes cai em $1$.
        \item O peso desse conjunto passa de $2^{-r}$ para $2^{-(r-1)} = 2 \times 2^{-r}$. O peso \textbf{dobra}.
        \item O aumento total no potencial gerado pelo Maker é exatamente a soma dos pesos locais de $x$:
        $$\Delta_{Maker} = \sum_{\substack{A \in \mathcal{F}, A \cap B_i = \emptyset \\ x \in A}} 2^{-|A \setminus M_i|}$$
    \end{itemize}
\end{frame}

\begin{frame}{Dinâmica das Jogadas: A Estratégia Gulosa do Breaker}
    \begin{itemize}
        \item Após a jogada do Maker, o potencial momentâneo é $\Phi_{i+0.5} = \Phi_i + \Delta_{Maker}$.
        \item Como $\Phi_0 < 1/2$ e a primeira jogada no máximo dobra esse valor, temos que $\Phi_{0.5} < 1$.
        \item \textbf{Ação do Breaker (Estratégia Gulosa):} O Breaker calcula o \textit{peso local} de todos os vértices livres $y$ e \textbf{escolhe aquele com o maior peso}.
        \item Ao escolher $y$, o Breaker "destrói" todos os conjuntos ativos que continham $y$, reduzindo o potencial do jogo em um montante $\Delta_{Breaker} = w(y)$.
        \item Como o Breaker escolhe o \textbf{máximo} disponível e o Maker já gastou sua jogada, matematicamente garante-se que:
        $$\Delta_{Breaker} \ge \Delta_{Maker}$$
        $$\implies \Phi_{i+1} \le \Phi_i$$
    \end{itemize}
\end{frame}

\begin{frame}{Conclusão da Demonstração por Contradição}
    \begin{block}{O Fechamento da Prova}
        \begin{itemize}
            \item Provamos que o potencial $\Phi$ nunca ultrapassa seu valor após a primeira jogada do Maker ($\Phi_{0.5} < 1$).
            \item Logo, o potencial \textbf{sempre permanecerá estritamente menor que 1}.
            \item Se o Maker vencesse na rodada $k$, ele teria ocupado inteiramente um conjunto $A^*$.
            \item O número de elementos faltantes em $A^*$ seria zero, e seu peso seria $2^{-0} = 1$.
            \item Isso forçaria o potencial total $\Phi_k$ a ser \textbf{pelo menos 1}.
            \item \textbf{Contradição!} Portanto, o Maker nunca poderá vencer. $\blacksquare$
        \end{itemize}
    \end{block}
\end{frame}

% ==========================================
% TÉCNICA B - DEFINIÇÃO E COTA INFERIOR
% ==========================================
\begin{frame}{Técnica B: Jogo Simétrico e Otimalidade do Critério}
    \small
    \textbf{Do Maker-Breaker ao Jogo Simétrico (\textit{Strong Game}):}
    \begin{itemize}
        \item O artigo original (1973) também considera o jogo em que ambos os jogadores têm o mesmo objetivo: \textbf{vence quem completar primeiro} algum conjunto $A \in \mathcal{F}$ (se ninguém completar, há empate).
        \item Seja $m^*(r)$ o menor número de conjuntos de tamanho $r$ que permite ao \textbf{Primeiro Jogador} forçar a vitória.
    \end{itemize}

    \begin{block}{Teorema da Construção Extremal (Erdős-Selfridge, 1973)}
        $$m^*(r) = 2^{r-1}$$
    \end{block}

    \textbf{Cota Inferior ($m^*(r) \ge 2^{r-1}$) via Técnica A:}
    \begin{itemize}
        \item Se $|\mathcal{F}| < 2^{r-1}$ e $|A| = r$ para todo $A \in \mathcal{F}$, então $\sum_{A \in \mathcal{F}} 2^{-|A|} = |\mathcal{F}| \cdot 2^{-r} < 2^{r-1} \cdot 2^{-r} = \tfrac{1}{2}$.
        \item Pela Técnica A, o 2º jogador bloqueia todos os conjuntos do 1º jogador (vencendo como \textit{Breaker} ou garantindo no mínimo o empate no jogo simétrico).
    \end{itemize}
\end{frame}

% ==========================================
% TÉCNICA B - CONSTRUÇÃO INDUTIVA
% ==========================================
\begin{frame}{Cota Superior ($m^*(r) \le 2^{r-1}$): Construção de Erdős-Selfridge}
    Para mostrar que $m^*(r) \le 2^{r-1}$, construímos por indução duas famílias gêmeas $G_r$ e $G'_r$, cada uma com $2^{r-1}$ conjuntos de tamanho $r$:

    \begin{itemize}
        \item \textbf{Base ($r=1$):} $G_1 = \{\{a_1\}\}$ e $G'_1 = \{\{b_1\}\}$ ($2^{1-1} = 1$ conjunto).
        \item \textbf{Passo Indutivo ($r-1 \to r$):} Dados $G_{r-1}$ e $G'_{r-1}$ (cada uma com $2^{r-2}$ conjuntos), adicionamos dois novos elementos $a_r$ e $b_r$:
        $$G_r = \{ A \cup \{a_r\} : A \in G_{r-1} \cup G'_{r-1} \}, \quad G'_r = \{ A \cup \{b_r\} : A \in G_{r-1} \cup G'_{r-1} \}$$
    \end{itemize}

    \begin{exampleblock}{Exemplos e Estrutura Explícita de $G_r$}
        \begin{itemize}
            \item $G_2 = \{\{a_1, a_2\}, \{b_1, a_2\}\}$ \quad ($2$ conjuntos de tamanho $2$)
            \item $G_3 = \{\{a_1, a_2, a_3\}, \{b_1, a_2, a_3\}, \{a_1, b_2, a_3\}, \{b_1, b_2, a_3\}\}$ \quad ($4$ conjuntos de tamanho $3$)
            \item \textbf{Regra geral:} Todo conjunto em $G_r$ contém o elemento fixo $\mathbf{a_r}$ mais \textbf{exatamente um elemento de cada par $\{a_i, b_i\}$} para $i \in \{1, \dots, r-1\}$.
        \end{itemize}
    \end{exampleblock}
\end{frame}

% ==========================================
% TÉCNICA B - ESTRATÉGIA VENCEDORA
% ==========================================
\begin{frame}{A Estratégia Vencedora (Emparelhamento) em $G_r$}
    Como o Primeiro Jogador força a vitória jogando sobre a família $G_r$?

    \begin{enumerate}
        \item \textbf{1º Lance (O golpe duplo em $a_r$):} O 1º jogador escolhe imediatamente $\mathbf{a_r}$.
        \begin{itemize}
            \item \textit{Defesa total no Jogo Simétrico:} Como $a_r \in A$ para todo $A \in G_r$, o 2º jogador fica \textbf{instantaneamente impedido de vencer}!
            \item \textit{Redução do ataque:} Resta ao 1º jogador conquistar um elemento de cada par $\{a_1, b_1\}, \dots, \{a_{r-1}, b_{r-1}\}$.
        \end{itemize}
        \vspace{0.2cm}
        \item \textbf{Rodadas $2$ até $r$ (Estratégia de Espelhamento):}
        \begin{itemize}
            \item Os únicos elementos restantes em $G_r$ são os $r-1$ pares disjuntos $\{a_i, b_i\}$.
            \item Sempre que o 2º jogador escolhe $a_i$ (bloqueando metade dos conjuntos ativos), o 1º jogador responde imediatamente escolhendo $b_i$ (e vice-versa).
        \end{itemize}
    \end{enumerate}

    \begin{alertblock}{Conclusão da Prova}
        No $r$-ésimo lance, o 1º jogador possui $a_r$ e um elemento $c_i \in \{a_i, b_i\}$ de cada par, completando $\{c_1, \dots, c_{r-1}, a_r\} \in G_r$ e vencendo tanto no Jogo Simétrico quanto no Maker-Breaker!
    \end{alertblock}
\end{frame}

\begin{frame}{Técnica C: Aplicação à Família de Cliques}
    Como aplicar isso para evitar a árvore de decisão e o "caos combinatório" de jogos em grafos?
    
    \begin{block}{O Jogo da Camarilha (Clique Game)}
        \begin{itemize}
            \item \textbf{Tabuleiro $X$:} Todas as arestas do grafo completo $K_N$ (Tamanho: $\binom{N}{2}$).
            \item \textbf{Objetivo do Maker:} Colorir arestas formando um clique $K_q$ completo.
            \item \textbf{Família $\mathcal{F}$:} Todos os possíveis cliques $K_q$ contidos em $K_N$.
        \end{itemize}
    \end{block}
    
    Quantos conjuntos vencedores existem? $\binom{N}{q}$.
    Qual o tamanho de cada conjunto vencedor? $\binom{q}{2}$.
\end{frame}

\begin{frame}{O Ponto de Quebra (Breaking Point)}
    Substituindo os valores no Critério de Erdős-Selfridge, o Breaker vence se:
    \begin{equation*}
        \sum_{A \in \mathcal{F}} 2^{-|A|} = \binom{N}{q} 2^{-\binom{q}{2}} < \frac{1}{2}
    \end{equation*}
    
    Para encontrar o limite assintótico:
    \begin{itemize}
        \item Aproximamos $\binom{N}{q} \approx \left(\frac{eN}{q}\right)^q \approx N^q$.
        \item A inequação vira $N^q \cdot 2^{-q^2/2} \approx 1/2$.
        \item Aplicando logaritmo na base 2: $q \log_2(N) - q^2/2 \approx -1$.
        \item Dividindo por $q$, obtemos a relação crítica:
    \end{itemize}
    \begin{alertblock}{O Limiar de Vitória}
        $q \approx 2 \log_2(N)$
    \end{alertblock}
    Acima desse tamanho de camarilha $q$, o Breaker sempre pode impedir o Maker usando puramente a estratégia do potencial!
\end{frame}

\begin{frame}{Conclusão}
    \begin{itemize}
        \item \textbf{Poder do Global sobre o Local:} O método do potencial avalia o estado geral do jogo de forma contínua, evitando o cálculo impossível de todas as respostas futuras.
        \item \textbf{O Nascimento do Método Quase-Probabilístico:} É uma ponte histórica entre a lógica de jogos determinísticos e a Teoria Probabilística.
        \item \textbf{Próximos Passos (Seminários 2 e 3):} O que acontece quando o jogo se torna \textbf{enviesado} ($a:b$)? Introduziremos os Box Games e os potenciais ponderados de Beck!
    \end{itemize}
    
    \vspace{1cm}
    \begin{center}
        \textbf{Obrigado!}
    \end{center}
\end{frame}

\end{document}